\documentclass[11pt]{article}
\usepackage[margin=2.2cm]{geometry}
\usepackage{amsmath,amssymb}
\usepackage{breqn}
\usepackage{needspace}
\usepackage[T1]{fontenc}
\begin{document}
\begin{center}
{\Large\bfseries TensorEngine: informe de corrida}
\end{center}
\textbf{Modelo:} \texttt{R}\\
\textbf{Run ID:} \texttt{run\_aa43c3ed782e2f3b0c98}\\
\textbf{Estado:} \texttt{partial}\\
\textbf{Backend:} \texttt{structural-python 0.9.0}\\
\textbf{Verificaciones:} 46 aprobadas, 0 fallidas, 2 indeterminadas.
\par\smallskip
\textbf{Presentación:} simplificación protegida, sin modificar los objetos canónicos. Las operaciones e hipótesis usadas por expresión constan en \texttt{presentation.json} y en los comentarios del archivo LaTeX. La validación xAct corresponde a la IR canónica.\par
\textbf{Deltas canónicos:} 819 sustituciones y 0 trazas registradas en las pasadas del álgebra (incluidas verificaciones); 0 deltas explícitos en las 13 cantidades finales. Detalle e índices sustituidos en \texttt{delta\_contractions.json}.\par
\section*{Expresiones tensoriales abstractas}
Las expresiones de esta sección son covariantes y no incorporan sustituciones del ansatz.
\Needspace{9\baselineskip}
\subsection*{$ L $}
% display: {"assumptions_used": [], "key": "abstract.lagrangian", "notes": [], "operations": [], "status": "unchanged"}
\begin{dmath*}[breakdepth={5}]
L = R{}_{d_{0}\,d_{1}\,d_{2}\,d_{3}}\,g{}^{d_{0}\,d_{2}}\,g{}^{d_{1}\,d_{3}}
\end{dmath*}
\par\vspace{1.2\baselineskip}\noindent
\begin{minipage}{\linewidth}
\textbf{Significado:} Densidad lagrangiana escalar normalizada de la teoría.\par
\textbf{Índices libres y varianzas:} $ \varnothing\;\text{(escalar)} $\par
\textbf{Fuentes del pipeline:} \texttt{validated\_model}\par
\textbf{Wolfram/xAct:} validación xAct no solicitada. La corrida no solicitó validación Wolfram/xAct.\par
\end{minipage}
\Needspace{9\baselineskip}
\subsection*{$ M_{ab} $}
% display: {"assumptions_used": [], "key": "abstract.metric_momentum", "notes": [], "operations": [], "status": "unchanged"}
\begin{dmath*}[breakdepth={5}]
M_{ab} = 2\,R{}_{a\,d_{0}\,b}{}^{d_{0}}
\end{dmath*}
\par\vspace{1.2\baselineskip}\noindent
\begin{minipage}{\linewidth}
\textbf{Significado:} Derivada de L respecto de la métrica inversa, M\_ab.\par
\textbf{Índices libres y varianzas:} $ a_{\downarrow},\;b_{\downarrow} $\par
\textbf{Fuentes del pipeline:} \texttt{lagrangian}\par
\textbf{Wolfram/xAct:} validación xAct no solicitada. La corrida no solicitó validación Wolfram/xAct.\par
\end{minipage}
\Needspace{9\baselineskip}
\subsection*{$ P^{abcd} $}
% display: {"assumptions_used": [], "key": "abstract.curvature_momentum", "notes": [], "operations": [], "status": "unchanged"}
\begin{dmath*}[breakdepth={5}]
P^{abcd} = -\frac{g{}^{a\,d}\,g{}^{b\,c}}{2} + \frac{g{}^{a\,c}\,g{}^{b\,d}}{2}
\end{dmath*}
\par\vspace{1.2\baselineskip}\noindent
\begin{minipage}{\linewidth}
\textbf{Significado:} Momento de curvatura P^{abcd}=partial L/partial R\_abcd.\par
\textbf{Índices libres y varianzas:} $ a^{\uparrow},\;b^{\uparrow},\;c^{\uparrow},\;d^{\uparrow} $\par
\textbf{Fuentes del pipeline:} \texttt{lagrangian}\par
\textbf{Wolfram/xAct:} validación xAct no solicitada. La corrida no solicitó validación Wolfram/xAct.\par
\end{minipage}
\Needspace{9\baselineskip}
\subsection*{$ J^a $}
% display: {"assumptions_used": [], "key": "abstract.scalar_gradient_momentum", "notes": [], "operations": [], "status": "unchanged"}
\begin{dmath*}[breakdepth={5}]
J^a = 0
\end{dmath*}
\par\vspace{1.2\baselineskip}\noindent
\begin{minipage}{\linewidth}
\textbf{Significado:} Momento J^a conjugado a nabla\_a phi.\par
\textbf{Índices libres y varianzas:} $ a^{\uparrow} $\par
\textbf{Fuentes del pipeline:} \texttt{lagrangian}\par
\textbf{Wolfram/xAct:} validación xAct no solicitada. La corrida no solicitó validación Wolfram/xAct.\par
\end{minipage}
\Needspace{9\baselineskip}
\subsection*{$ F_{\phi} $}
% display: {"assumptions_used": [], "key": "abstract.scalar_derivative", "notes": [], "operations": [], "status": "unchanged"}
\begin{dmath*}[breakdepth={5}]
F_{\phi} = 0
\end{dmath*}
\par\vspace{1.2\baselineskip}\noindent
\begin{minipage}{\linewidth}
\textbf{Significado:} Derivada explícita F\_phi=partial L/partial phi.\par
\textbf{Índices libres y varianzas:} $ \varnothing\;\text{(escalar)} $\par
\textbf{Fuentes del pipeline:} \texttt{lagrangian}\par
\textbf{Wolfram/xAct:} validación xAct no solicitada. La corrida no solicitó validación Wolfram/xAct.\par
\end{minipage}
\Needspace{9\baselineskip}
\subsection*{$ E_{ab} $}
% display: {"assumptions_used": [], "key": "abstract.metric_euler", "notes": [], "operations": [], "status": "unchanged"}
\begin{dmath*}[breakdepth={5}]
E_{ab} = -\frac{g{}_{a\,b}\,R{}_{d_{0}\,d_{1}}{}^{d_{0}\,d_{1}}}{2} + R{}_{a\,d_{0}\,b}{}^{d_{0}}
\end{dmath*}
\par\vspace{1.2\baselineskip}\noindent
\begin{minipage}{\linewidth}
\textbf{Significado:} Ecuación de Euler-Lagrange métrica E\_ab.\par
\textbf{Índices libres y varianzas:} $ a_{\downarrow},\;b_{\downarrow} $\par
\textbf{Fuentes del pipeline:} \texttt{delta\_lagrangian, curvature\_derivative\_metric\_term}\par
\textbf{Wolfram/xAct:} validación xAct no solicitada. La corrida no solicitó validación Wolfram/xAct.\par
\end{minipage}
% display: {"assumptions_used": [], "key": "abstract.curvature_derivative_metric_term", "notes": [], "operations": [], "status": "unchanged"}
\begin{dmath*}[breakdepth={5}]
\left.E_{ab}\right|_{\nabla\nabla P} = 0
\end{dmath*}
\par\vspace{0.8\baselineskip}\noindent
\begin{minipage}{\linewidth}
La expresión anterior identifica la contribución $-2\nabla^c\nabla^dP_{acdb}$ dentro de $E_{ab}$.\par
\end{minipage}
\Needspace{9\baselineskip}
\subsection*{$ E_{\phi} $}
% display: {"assumptions_used": [], "key": "abstract.scalar_euler", "notes": [], "operations": [], "status": "unchanged"}
\begin{dmath*}[breakdepth={5}]
E_{\phi} = 0
\end{dmath*}
\par\vspace{1.2\baselineskip}\noindent
\begin{minipage}{\linewidth}
\textbf{Significado:} Ecuación de Euler-Lagrange escalar E\_phi.\par
\textbf{Índices libres y varianzas:} $ \varnothing\;\text{(escalar)} $\par
\textbf{Fuentes del pipeline:} \texttt{scalar\_derivative, scalar\_gradient\_momentum}\par
\textbf{Wolfram/xAct:} validación xAct no solicitada. La corrida no solicitó validación Wolfram/xAct.\par
\end{minipage}
\Needspace{9\baselineskip}
\subsection*{$ \mathcal{R} $}
% display: {"assumptions_used": [], "key": "abstract.ricci_scalar", "notes": [], "operations": [], "status": "unchanged"}
\begin{dmath*}[breakdepth={5}]
\mathcal{R} = g{}^{a\,c}\,g{}^{b\,d}\,R{}_{a\,b\,c\,d}
\end{dmath*}
\par\vspace{1.2\baselineskip}\noindent
\begin{minipage}{\linewidth}
\textbf{Significado:} Escalar de Ricci construido con la convención geométrica activa.\par
\textbf{Índices libres y varianzas:} $ \varnothing\;\text{(escalar)} $\par
\textbf{Fuentes del pipeline:} \texttt{validated\_model}\par
\textbf{Wolfram/xAct:} validación xAct no solicitada. La corrida no solicitó validación Wolfram/xAct.\par
\end{minipage}
\Needspace{9\baselineskip}
\subsection*{$ R_{ab}R^{ab} $}
% display: {"assumptions_used": [], "key": "abstract.ricci_squared", "notes": [], "operations": [], "status": "unchanged"}
\begin{dmath*}[breakdepth={5}]
R_{ab}R^{ab} = g{}^{a\,c}\,R{}_{a\,b\,c\,d}\,g{}^{b\,e}\,g{}^{d\,f}\,g{}^{g\,h}\,R{}_{g\,e\,h\,f}
\end{dmath*}
\par\vspace{1.2\baselineskip}\noindent
\begin{minipage}{\linewidth}
\textbf{Significado:} Invariante cuadrático de Ricci R\_ab R^ab.\par
\textbf{Índices libres y varianzas:} $ \varnothing\;\text{(escalar)} $\par
\textbf{Fuentes del pipeline:} \texttt{validated\_model, riemann\_tensor}\par
\textbf{Wolfram/xAct:} validación xAct no solicitada. La corrida no solicitó validación Wolfram/xAct.\par
\end{minipage}
\Needspace{9\baselineskip}
\subsection*{$ R_{abcd} $}
% display: {"assumptions_used": [], "key": "abstract.riemann_tensor", "notes": [], "operations": [], "status": "unchanged"}
\begin{dmath*}[breakdepth={5}]
R_{abcd} = R{}_{a\,b\,c\,d}
\end{dmath*}
\par\vspace{1.2\baselineskip}\noindent
\begin{minipage}{\linewidth}
\textbf{Significado:} Tensor de Riemann covariante tratado como entrada geométrica.\par
\textbf{Índices libres y varianzas:} $ a_{\downarrow},\;b_{\downarrow},\;c_{\downarrow},\;d_{\downarrow} $\par
\textbf{Fuentes del pipeline:} \texttt{validated\_model}\par
\textbf{Wolfram/xAct:} validación xAct no solicitada. La corrida no solicitó validación Wolfram/xAct.\par
\end{minipage}
\Needspace{9\baselineskip}
\subsection*{$ R_{abcd}R^{abcd} $}
% display: {"assumptions_used": [], "key": "abstract.riemann_squared", "notes": [], "operations": [], "status": "unchanged"}
\begin{dmath*}[breakdepth={5}]
R_{abcd}R^{abcd} = g{}^{a\,e}\,g{}^{b\,f}\,g{}^{c\,g}\,g{}^{d\,h}\,R{}_{a\,b\,c\,d}\,R{}_{e\,f\,g\,h}
\end{dmath*}
\par\vspace{1.2\baselineskip}\noindent
\begin{minipage}{\linewidth}
\textbf{Significado:} Invariante de Kretschmann R\_abcd R^abcd.\par
\textbf{Índices libres y varianzas:} $ \varnothing\;\text{(escalar)} $\par
\textbf{Fuentes del pipeline:} \texttt{validated\_model, riemann\_tensor}\par
\textbf{Wolfram/xAct:} validación xAct no solicitada. La corrida no solicitó validación Wolfram/xAct.\par
\end{minipage}
\Needspace{9\baselineskip}
\subsection*{$ \nabla_e P^{abcd} $}
% display: {"assumptions_used": [], "key": "abstract.nabla_P", "notes": [], "operations": [], "status": "unchanged"}
\begin{dmath*}[breakdepth={5}]
\nabla_e P^{abcd} = 0
\end{dmath*}
\par\vspace{1.2\baselineskip}\noindent
\begin{minipage}{\linewidth}
\textbf{Significado:} Primera derivada covariante del momento de curvatura.\par
\textbf{Índices libres y varianzas:} $ a^{\uparrow},\;b^{\uparrow},\;c^{\uparrow},\;d^{\uparrow},\;e_{\downarrow} $\par
\textbf{Fuentes del pipeline:} \texttt{curvature\_momentum}\par
\textbf{Wolfram/xAct:} validación xAct no solicitada. La corrida no solicitó validación Wolfram/xAct.\par
\end{minipage}
\Needspace{9\baselineskip}
\subsection*{$ \nabla_f\nabla_e P^{abcd} $}
% display: {"assumptions_used": [], "key": "abstract.nabla_nabla_P", "notes": [], "operations": [], "status": "unchanged"}
\begin{dmath*}[breakdepth={5}]
\nabla_f\nabla_e P^{abcd} = 0
\end{dmath*}
\par\vspace{1.2\baselineskip}\noindent
\begin{minipage}{\linewidth}
\textbf{Significado:} Segunda derivada covariante del momento de curvatura.\par
\textbf{Índices libres y varianzas:} $ a^{\uparrow},\;b^{\uparrow},\;c^{\uparrow},\;d^{\uparrow},\;e_{\downarrow},\;f_{\downarrow} $\par
\textbf{Fuentes del pipeline:} \texttt{nabla\_P}\par
\textbf{Wolfram/xAct:} validación xAct no solicitada. La corrida no solicitó validación Wolfram/xAct.\par
\end{minipage}
\Needspace{12\baselineskip}
\par\bigskip\noindent{\large\bfseries Descomposición compacta adicional}\par
Esta vista se agrega después de los resultados anteriores y reutiliza exactamente la IR canónica ya calculada. No interviene en la validación ni en los fingerprints.\par
\Needspace{10\baselineskip}
\par\medskip\noindent\textbf{Descomposición de $ E_{ab} $.}\par
\begin{dmath*}[breakdepth={5}]
E_{ab}=M_{ab}-\mathcal{R}_{(ab)}[P]-\frac{1}{2}g_{ab}L-2\nabla_m\nabla_nP^{mn}{}_{ab}
\end{dmath*}
\textbf{Reconstrucción IR:} verified. Los cuatro bloques reconstruyen exactamente E\_ab en la IR canónica.\par
% compact-display: {"assumptions_used": [], "key": "abstract.compact.metric_euler", "notes": [], "operations": [], "status": "unchanged"}
\begin{dmath*}[breakdepth={5}]
\text{forma compacta:}\quad E_{ab} = -\frac{g{}_{a\,b}\,R{}_{d_{0}\,d_{1}}{}^{d_{0}\,d_{1}}}{2} + R{}_{a\,d_{0}\,b}{}^{d_{0}}
\end{dmath*}
\par\smallskip\noindent\textbf{Bloque $ M_{ab} $} (fuentes: \texttt{metric\_momentum}).\par
% compact-display: {"assumptions_used": [], "key": "abstract.compact.metric_euler.metric_momentum", "notes": [], "operations": [], "status": "unchanged"}
\begin{dmath*}[breakdepth={5}]
\text{forma compacta:}\quad M_{ab} = 2\,R{}_{a\,d_{0}\,b}{}^{d_{0}}
\end{dmath*}
\par\smallskip\noindent\textbf{Bloque $ -\mathcal{R}_{(ab)}[P] $} (fuentes: \texttt{curvature\_momentum, riemann\_tensor}).\par
% compact-display: {"assumptions_used": [], "key": "abstract.compact.metric_euler.curvature_algebraic", "notes": [], "operations": [], "status": "unchanged"}
\begin{dmath*}[breakdepth={5}]
\text{forma compacta:}\quad -\mathcal{R}_{(ab)}[P] = \frac{1}{2}\,\left(-g{}_{a\,d_{0}}\,\left(-\frac{g{}^{d_{0}\,d_{1}}\,g{}^{d_{2}\,d_{3}}}{2} + \frac{g{}^{d_{0}\,d_{3}}\,g{}^{d_{2}\,d_{1}}}{2}\right)\,R{}_{b\,d_{2}\,d_{3}\,d_{1}} - g{}_{b\,d_{0}}\,\left(-\frac{g{}^{d_{0}\,d_{1}}\,g{}^{d_{2}\,d_{3}}}{2} + \frac{g{}^{d_{0}\,d_{3}}\,g{}^{d_{2}\,d_{1}}}{2}\right)\,R{}_{a\,d_{2}\,d_{3}\,d_{1}}\right)
\end{dmath*}
\par\smallskip\noindent\textbf{Bloque $ -\frac{1}{2}g_{ab}L $} (fuentes: \texttt{lagrangian}).\par
% compact-display: {"assumptions_used": [], "key": "abstract.compact.metric_euler.volume_term", "notes": [], "operations": [], "status": "unchanged"}
\begin{dmath*}[breakdepth={5}]
\text{forma compacta:}\quad -\frac{1}{2}g_{ab}L = -\frac{g{}_{a\,b}\,R{}_{d_{0}\,d_{1}\,d_{2}\,d_{3}}\,g{}^{d_{0}\,d_{2}}\,g{}^{d_{1}\,d_{3}}}{2}
\end{dmath*}
\par\smallskip\noindent\textbf{Bloque $ E_{ab}^{(\nabla\nabla P)} $} (fuentes: \texttt{curvature\_momentum, nabla\_nabla\_P}).\par
% compact-display: {"assumptions_used": [], "key": "abstract.compact.metric_euler.curvature_derivative", "notes": [], "operations": [], "status": "unchanged"}
\begin{dmath*}[breakdepth={5}]
\text{forma compacta:}\quad E_{ab}^{(\nabla\nabla P)} = 0
\end{dmath*}
\Needspace{10\baselineskip}
\par\medskip\noindent\textbf{Descomposición de $ E_{\phi} $.}\par
\begin{dmath*}[breakdepth={5}]
E_{\phi}=F_{\phi}-\nabla_aJ^a
\end{dmath*}
\textbf{Reconstrucción IR:} verified. Los dos bloques reconstruyen exactamente E\_phi en la IR canónica.\par
% compact-display: {"assumptions_used": [], "key": "abstract.compact.scalar_euler", "notes": [], "operations": [], "status": "unchanged"}
\begin{dmath*}[breakdepth={5}]
\text{forma compacta:}\quad E_{\phi} = 0
\end{dmath*}
\par\smallskip\noindent\textbf{Bloque $ F_{\phi} $} (fuentes: \texttt{scalar\_derivative}).\par
% compact-display: {"assumptions_used": [], "key": "abstract.compact.scalar_euler.scalar_force", "notes": [], "operations": [], "status": "unchanged"}
\begin{dmath*}[breakdepth={5}]
\text{forma compacta:}\quad F_{\phi} = 0
\end{dmath*}
\par\smallskip\noindent\textbf{Bloque $ -\nabla_a J^a $} (fuentes: \texttt{scalar\_gradient\_momentum, scalar\_euler}).\par
% compact-display: {"assumptions_used": [], "key": "abstract.compact.scalar_euler.scalar_current_divergence", "notes": [], "operations": [], "status": "unchanged"}
\begin{dmath*}[breakdepth={5}]
\text{forma compacta:}\quad -\nabla_a J^a = 0
\end{dmath*}
\Needspace{10\baselineskip}
\par\medskip\noindent\textbf{Descomposición de $ P^{abcd} $.}\par
\begin{dmath*}[breakdepth={5}]
P^{abcd}=\partial L/\partial R_{abcd}
\end{dmath*}
\textbf{Reconstrucción IR:} verified. La forma compacta y la expandida referencian el mismo objeto P^{abcd}.\par
% compact-display: {"assumptions_used": [], "key": "abstract.compact.curvature_momentum", "notes": [], "operations": [], "status": "unchanged"}
\begin{dmath*}[breakdepth={5}]
\text{forma compacta:}\quad P^{abcd} = -\frac{g{}^{a\,d}\,g{}^{b\,c}}{2} + \frac{g{}^{a\,c}\,g{}^{b\,d}}{2}
\end{dmath*}
\section*{Expresiones proyectadas mediante el ansatz}
\textbf{Ansatz utilizado:} \texttt{draft4\_circular}.
\begin{dmath*}[breakdepth={3}]
(x^\mu)=\left(\tau, r, \varphi\right),\qquad ds^2 = -f\!\left(r\right)\,d\tau^2 + \frac{1}{f\!\left(r\right)}\,dr^2 + {r}^{2}\,d\varphi^2
\end{dmath*}
\begin{dmath*}[breakdepth={3}]
\phi = \Phi\!\left(r, \varphi\right)\qquad\text{(genérico, sin especialización)}
\end{dmath*}
Las componentes completas se conservan de forma dispersa en \texttt{results.json}; el documento muestra como máximo doce componentes no nulas por tensor.
\Needspace{7\baselineskip}
\subsection*{$ L $}
\textbf{Ansatz:} \texttt{draft4\_circular}; \textbf{estado:} completada.\par
% display: {"assumptions_used": [], "key": "projected.lagrangian.scalar", "notes": [], "operations": [], "status": "unchanged"}
\begin{dmath*}[breakdepth={5}]
L\big|_{\mathrm{ansatz}} = -\left(\frac{2\,f'\!\left(r\right)}{r} + f''\!\left(r\right)\right)
\end{dmath*}
\par\smallskip
Proyección completa almacenada: 1 componentes no nulas de 1. Proyectada respetando el ansatz 'draft4\_circular'.\par\medskip
\Needspace{7\baselineskip}
\subsection*{$ M_{ab} $}
\textbf{Ansatz:} \texttt{draft4\_circular}; \textbf{estado:} completada.\par
% display: {"assumptions_used": [], "key": "projected.metric_momentum[0,0]", "notes": [], "operations": [], "status": "unchanged"}
\begin{dmath*}[breakdepth={5}]
\left[M_{ab}\right]_{00} = \frac{\left(r\,f''\!\left(r\right) + f'\!\left(r\right)\right)\,f\!\left(r\right)}{r}
\end{dmath*}
\par\smallskip
% display: {"assumptions_used": [], "key": "projected.metric_momentum[1,1]", "notes": [], "operations": [], "status": "unchanged"}
\begin{dmath*}[breakdepth={5}]
\left[M_{ab}\right]_{11} = -\frac{\left(r\,f''\!\left(r\right) + f'\!\left(r\right)\right)}{r\,f\!\left(r\right)}
\end{dmath*}
\par\smallskip
% display: {"assumptions_used": [], "key": "projected.metric_momentum[2,2]", "notes": [], "operations": [], "status": "unchanged"}
\begin{dmath*}[breakdepth={5}]
\left[M_{ab}\right]_{22} = -2\,r\,f'\!\left(r\right)
\end{dmath*}
\par\smallskip
Proyección completa almacenada: 3 componentes no nulas de 9. Proyectada respetando el ansatz 'draft4\_circular'.\par\medskip
\Needspace{7\baselineskip}
\subsection*{$ P^{abcd} $}
\textbf{Ansatz:} \texttt{draft4\_circular}; \textbf{estado:} completada.\par
% display: {"assumptions_used": [], "key": "projected.curvature_momentum[0,0,1,1]", "notes": [], "operations": [], "status": "unchanged"}
\begin{dmath*}[breakdepth={5}]
\left[P^{abcd}\right]^{0110} = \frac{1}{2}
\end{dmath*}
\par\smallskip
% display: {"assumptions_used": [], "key": "projected.curvature_momentum[0,0,2,2]", "notes": [], "operations": [], "status": "unchanged"}
\begin{dmath*}[breakdepth={5}]
\left[P^{abcd}\right]^{0220} = \frac{1}{2\,{r}^{2}\,f\!\left(r\right)}
\end{dmath*}
\par\smallskip
% display: {"assumptions_used": [], "key": "projected.curvature_momentum[0,1,1,0]", "notes": [], "operations": [], "status": "unchanged"}
\begin{dmath*}[breakdepth={5}]
\left[P^{abcd}\right]^{0101} = -\frac{1}{2}
\end{dmath*}
\par\smallskip
% display: {"assumptions_used": [], "key": "projected.curvature_momentum[0,2,2,0]", "notes": [], "operations": [], "status": "unchanged"}
\begin{dmath*}[breakdepth={5}]
\left[P^{abcd}\right]^{0202} = -\frac{1}{2\,{r}^{2}\,f\!\left(r\right)}
\end{dmath*}
\par\smallskip
% display: {"assumptions_used": [], "key": "projected.curvature_momentum[1,0,0,1]", "notes": [], "operations": [], "status": "unchanged"}
\begin{dmath*}[breakdepth={5}]
\left[P^{abcd}\right]^{1010} = -\frac{1}{2}
\end{dmath*}
\par\smallskip
% display: {"assumptions_used": [], "key": "projected.curvature_momentum[1,1,0,0]", "notes": [], "operations": [], "status": "unchanged"}
\begin{dmath*}[breakdepth={5}]
\left[P^{abcd}\right]^{1001} = \frac{1}{2}
\end{dmath*}
\par\smallskip
% display: {"assumptions_used": [], "key": "projected.curvature_momentum[1,1,2,2]", "notes": [], "operations": [], "status": "unchanged"}
\begin{dmath*}[breakdepth={5}]
\left[P^{abcd}\right]^{1221} = -\frac{f\!\left(r\right)}{2\,{r}^{2}}
\end{dmath*}
\par\smallskip
% display: {"assumptions_used": [], "key": "projected.curvature_momentum[1,2,2,1]", "notes": [], "operations": [], "status": "unchanged"}
\begin{dmath*}[breakdepth={5}]
\left[P^{abcd}\right]^{1212} = \frac{f\!\left(r\right)}{2\,{r}^{2}}
\end{dmath*}
\par\smallskip
% display: {"assumptions_used": [], "key": "projected.curvature_momentum[2,0,0,2]", "notes": [], "operations": [], "status": "unchanged"}
\begin{dmath*}[breakdepth={5}]
\left[P^{abcd}\right]^{2020} = -\frac{1}{2\,{r}^{2}\,f\!\left(r\right)}
\end{dmath*}
\par\smallskip
% display: {"assumptions_used": [], "key": "projected.curvature_momentum[2,1,1,2]", "notes": [], "operations": [], "status": "unchanged"}
\begin{dmath*}[breakdepth={5}]
\left[P^{abcd}\right]^{2121} = \frac{f\!\left(r\right)}{2\,{r}^{2}}
\end{dmath*}
\par\smallskip
% display: {"assumptions_used": [], "key": "projected.curvature_momentum[2,2,0,0]", "notes": [], "operations": [], "status": "unchanged"}
\begin{dmath*}[breakdepth={5}]
\left[P^{abcd}\right]^{2002} = \frac{1}{2\,{r}^{2}\,f\!\left(r\right)}
\end{dmath*}
\par\smallskip
% display: {"assumptions_used": [], "key": "projected.curvature_momentum[2,2,1,1]", "notes": [], "operations": [], "status": "unchanged"}
\begin{dmath*}[breakdepth={5}]
\left[P^{abcd}\right]^{2112} = -\frac{f\!\left(r\right)}{2\,{r}^{2}}
\end{dmath*}
\par\smallskip
Proyección completa almacenada: 12 componentes no nulas de 81. Proyectada respetando el ansatz 'draft4\_circular'.\par\medskip
\Needspace{7\baselineskip}
\subsection*{$ J^a $}
\textbf{Ansatz:} \texttt{draft4\_circular}; \textbf{estado:} completada.\par
% display: {"assumptions_used": [], "key": "projected.scalar_gradient_momentum.zero", "notes": [], "operations": [], "status": "unchanged"}
Todas las componentes de esta cantidad son nulas.
Proyección completa almacenada: 0 componentes no nulas de 3. La expresión abstracta es nula; todas sus componentes son cero.\par\medskip
\Needspace{7\baselineskip}
\subsection*{$ F_{\phi} $}
\textbf{Ansatz:} \texttt{draft4\_circular}; \textbf{estado:} completada.\par
% display: {"assumptions_used": [], "key": "projected.scalar_derivative.scalar", "notes": [], "operations": [], "status": "unchanged"}
\begin{dmath*}[breakdepth={5}]
F_{\phi}\big|_{\mathrm{ansatz}} = 0
\end{dmath*}
\par\smallskip
Proyección completa almacenada: 0 componentes no nulas de 1. La expresión abstracta es nula; todas sus componentes son cero.\par\medskip
\Needspace{7\baselineskip}
\subsection*{$ E_{ab} $}
\textbf{Ansatz:} \texttt{draft4\_circular}; \textbf{estado:} completada.\par
% display: {"assumptions_used": [], "key": "projected.metric_euler[0,0]", "notes": [], "operations": [], "status": "unchanged"}
\begin{dmath*}[breakdepth={5}]
\left[E_{ab}\right]_{00} = -\frac{f'\!\left(r\right)\,f\!\left(r\right)}{2\,r}
\end{dmath*}
\par\smallskip
% display: {"assumptions_used": [], "key": "projected.metric_euler[1,1]", "notes": [], "operations": [], "status": "unchanged"}
\begin{dmath*}[breakdepth={5}]
\left[E_{ab}\right]_{11} = \frac{f'\!\left(r\right)}{2\,r\,f\!\left(r\right)}
\end{dmath*}
\par\smallskip
% display: {"assumptions_used": [], "key": "projected.metric_euler[2,2]", "notes": [], "operations": [], "status": "unchanged"}
\begin{dmath*}[breakdepth={5}]
\left[E_{ab}\right]_{22} = \frac{{r}^{2}\,f''\!\left(r\right)}{2}
\end{dmath*}
\par\smallskip
Proyección completa almacenada: 3 componentes no nulas de 9. Reutilizada desde la etapa de componentes de E\_ab.\par\medskip
\Needspace{7\baselineskip}
\subsection*{$ E_{\phi} $}
\textbf{Ansatz:} \texttt{draft4\_circular}; \textbf{estado:} completada.\par
% display: {"assumptions_used": [], "key": "projected.scalar_euler.scalar", "notes": [], "operations": [], "status": "unchanged"}
\begin{dmath*}[breakdepth={5}]
E_{\phi}\big|_{\mathrm{ansatz}} = 0
\end{dmath*}
\par\smallskip
Proyección completa almacenada: 0 componentes no nulas de 1. Reutilizada desde la etapa de componentes de E\_phi.\par\medskip
\Needspace{7\baselineskip}
\subsection*{$ \mathcal{R} $}
\textbf{Ansatz:} \texttt{draft4\_circular}; \textbf{estado:} completada.\par
% display: {"assumptions_used": [], "key": "projected.ricci_scalar.scalar", "notes": [], "operations": [], "status": "unchanged"}
\begin{dmath*}[breakdepth={5}]
\mathcal{R}\big|_{\mathrm{ansatz}} = -\left(\frac{2\,f'\!\left(r\right)}{r} + f''\!\left(r\right)\right)
\end{dmath*}
\par\smallskip
Proyección completa almacenada: 1 componentes no nulas de 1. Proyectada respetando el ansatz 'draft4\_circular'. La corrida no solicitó validación Wolfram/xAct.\par\medskip
\Needspace{7\baselineskip}
\subsection*{$ R_{ab}R^{ab} $}
\textbf{Ansatz:} \texttt{draft4\_circular}; \textbf{estado:} completada.\par
% display: {"assumptions_used": [], "key": "projected.ricci_squared.scalar", "notes": [], "operations": [], "status": "unchanged"}
\begin{dmath*}[breakdepth={5}]
R_{ab}R^{ab}\big|_{\mathrm{ansatz}} = \frac{{f''\!\left(r\right)}^{2}}{2} + \frac{3\,{f'\!\left(r\right)}^{2}}{2\,{r}^{2}} + \frac{f'\!\left(r\right)\,f''\!\left(r\right)}{r}
\end{dmath*}
\par\smallskip
Proyección completa almacenada: 1 componentes no nulas de 1. Proyectada respetando el ansatz 'draft4\_circular'. La corrida no solicitó validación Wolfram/xAct.\par\medskip
\Needspace{7\baselineskip}
\subsection*{$ R_{abcd} $}
\textbf{Ansatz:} \texttt{draft4\_circular}; \textbf{estado:} completada.\par
% display: {"assumptions_used": [], "key": "projected.riemann_tensor[0,1,0,1]", "notes": [], "operations": [], "status": "unchanged"}
\begin{dmath*}[breakdepth={5}]
\left[R_{abcd}\right]_{0101} = \frac{f''\!\left(r\right)}{2}
\end{dmath*}
\par\smallskip
% display: {"assumptions_used": [], "key": "projected.riemann_tensor[0,1,1,0]", "notes": [], "operations": [], "status": "unchanged"}
\begin{dmath*}[breakdepth={5}]
\left[R_{abcd}\right]_{0110} = -\frac{f''\!\left(r\right)}{2}
\end{dmath*}
\par\smallskip
% display: {"assumptions_used": [], "key": "projected.riemann_tensor[0,2,0,2]", "notes": [], "operations": [], "status": "unchanged"}
\begin{dmath*}[breakdepth={5}]
\left[R_{abcd}\right]_{0202} = \frac{r\,f'\!\left(r\right)\,f\!\left(r\right)}{2}
\end{dmath*}
\par\smallskip
% display: {"assumptions_used": [], "key": "projected.riemann_tensor[0,2,2,0]", "notes": [], "operations": [], "status": "unchanged"}
\begin{dmath*}[breakdepth={5}]
\left[R_{abcd}\right]_{0220} = -\frac{r\,f'\!\left(r\right)\,f\!\left(r\right)}{2}
\end{dmath*}
\par\smallskip
% display: {"assumptions_used": [], "key": "projected.riemann_tensor[1,0,0,1]", "notes": [], "operations": [], "status": "unchanged"}
\begin{dmath*}[breakdepth={5}]
\left[R_{abcd}\right]_{1001} = -\frac{f''\!\left(r\right)}{2}
\end{dmath*}
\par\smallskip
% display: {"assumptions_used": [], "key": "projected.riemann_tensor[1,0,1,0]", "notes": [], "operations": [], "status": "unchanged"}
\begin{dmath*}[breakdepth={5}]
\left[R_{abcd}\right]_{1010} = \frac{f''\!\left(r\right)}{2}
\end{dmath*}
\par\smallskip
% display: {"assumptions_used": [], "key": "projected.riemann_tensor[1,2,1,2]", "notes": [], "operations": [], "status": "unchanged"}
\begin{dmath*}[breakdepth={5}]
\left[R_{abcd}\right]_{1212} = -\frac{r\,f'\!\left(r\right)}{2\,f\!\left(r\right)}
\end{dmath*}
\par\smallskip
% display: {"assumptions_used": [], "key": "projected.riemann_tensor[1,2,2,1]", "notes": [], "operations": [], "status": "unchanged"}
\begin{dmath*}[breakdepth={5}]
\left[R_{abcd}\right]_{1221} = \frac{r\,f'\!\left(r\right)}{2\,f\!\left(r\right)}
\end{dmath*}
\par\smallskip
% display: {"assumptions_used": [], "key": "projected.riemann_tensor[2,0,0,2]", "notes": [], "operations": [], "status": "unchanged"}
\begin{dmath*}[breakdepth={5}]
\left[R_{abcd}\right]_{2002} = -\frac{r\,f'\!\left(r\right)\,f\!\left(r\right)}{2}
\end{dmath*}
\par\smallskip
% display: {"assumptions_used": [], "key": "projected.riemann_tensor[2,0,2,0]", "notes": [], "operations": [], "status": "unchanged"}
\begin{dmath*}[breakdepth={5}]
\left[R_{abcd}\right]_{2020} = \frac{r\,f'\!\left(r\right)\,f\!\left(r\right)}{2}
\end{dmath*}
\par\smallskip
% display: {"assumptions_used": [], "key": "projected.riemann_tensor[2,1,1,2]", "notes": [], "operations": [], "status": "unchanged"}
\begin{dmath*}[breakdepth={5}]
\left[R_{abcd}\right]_{2112} = \frac{r\,f'\!\left(r\right)}{2\,f\!\left(r\right)}
\end{dmath*}
\par\smallskip
% display: {"assumptions_used": [], "key": "projected.riemann_tensor[2,1,2,1]", "notes": [], "operations": [], "status": "unchanged"}
\begin{dmath*}[breakdepth={5}]
\left[R_{abcd}\right]_{2121} = -\frac{r\,f'\!\left(r\right)}{2\,f\!\left(r\right)}
\end{dmath*}
\par\smallskip
Proyección completa almacenada: 12 componentes no nulas de 81. Proyectada respetando el ansatz 'draft4\_circular'. La corrida no solicitó validación Wolfram/xAct.\par\medskip
\Needspace{7\baselineskip}
\subsection*{$ R_{abcd}R^{abcd} $}
\textbf{Ansatz:} \texttt{draft4\_circular}; \textbf{estado:} completada.\par
% display: {"assumptions_used": [], "key": "projected.riemann_squared.scalar", "notes": [], "operations": [], "status": "unchanged"}
\begin{dmath*}[breakdepth={5}]
R_{abcd}R^{abcd}\big|_{\mathrm{ansatz}} = {f''\!\left(r\right)}^{2} + \frac{2\,{f'\!\left(r\right)}^{2}}{{r}^{2}}
\end{dmath*}
\par\smallskip
Proyección completa almacenada: 1 componentes no nulas de 1. Proyectada respetando el ansatz 'draft4\_circular'. La corrida no solicitó validación Wolfram/xAct.\par\medskip
\Needspace{7\baselineskip}
\subsection*{$ \nabla_e P^{abcd} $}
\textbf{Ansatz:} \texttt{draft4\_circular}; \textbf{estado:} completada.\par
% display: {"assumptions_used": [], "key": "projected.nabla_P.zero", "notes": [], "operations": [], "status": "unchanged"}
Todas las componentes de esta cantidad son nulas.
Proyección completa almacenada: 0 componentes no nulas de 243. La expresión abstracta es nula; todas sus componentes son cero. La corrida no solicitó validación Wolfram/xAct.\par\medskip
\Needspace{7\baselineskip}
\subsection*{$ \nabla_f\nabla_e P^{abcd} $}
\textbf{Ansatz:} \texttt{draft4\_circular}; \textbf{estado:} completada.\par
% display: {"assumptions_used": [], "key": "projected.nabla_nabla_P.zero", "notes": [], "operations": [], "status": "unchanged"}
Todas las componentes de esta cantidad son nulas.
Proyección completa almacenada: 0 componentes no nulas de 729. La expresión abstracta es nula; todas sus componentes son cero. La corrida no solicitó validación Wolfram/xAct.\par\medskip
\par\medskip Extras: perfiles escalares de una variable, con la misma métrica del ansatz genérico. Se sustituyen únicamente las componentes de $L$ y $P^{abcd}$ ya calculadas; sin validación xAct independiente de estos perfiles. Las componentes completas de estos extras constan en \texttt{presentation.json}.
\Needspace{8\baselineskip}
\subsection*{Extra: $\phi=\Phi\!\left(r\right)$}
\begin{dmath*}[breakdepth={5}]
L=-\left(\frac{2\,f'\!\left(r\right)}{r} + f''\!\left(r\right)\right)
\end{dmath*}
Componentes independientes no nulas; $P^{abcd}=-P^{bacd}=-P^{abdc}=P^{cdab}$.
\begin{dmath*}[breakdepth={5}]
\left[P^{abcd}\right]^{0101}=-\frac{1}{2}
\end{dmath*}
\begin{dmath*}[breakdepth={5}]
\left[P^{abcd}\right]^{0202}=-\frac{1}{2\,{r}^{2}\,f\!\left(r\right)}
\end{dmath*}
\begin{dmath*}[breakdepth={5}]
\left[P^{abcd}\right]^{1212}=\frac{f\!\left(r\right)}{2\,{r}^{2}}
\end{dmath*}
\Needspace{8\baselineskip}
\subsection*{Extra: $\phi=\Phi\!\left(\varphi\right)$}
\begin{dmath*}[breakdepth={5}]
L=-\left(\frac{2\,f'\!\left(r\right)}{r} + f''\!\left(r\right)\right)
\end{dmath*}
Componentes independientes no nulas; $P^{abcd}=-P^{bacd}=-P^{abdc}=P^{cdab}$.
\begin{dmath*}[breakdepth={5}]
\left[P^{abcd}\right]^{0101}=-\frac{1}{2}
\end{dmath*}
\begin{dmath*}[breakdepth={5}]
\left[P^{abcd}\right]^{0202}=-\frac{1}{2\,{r}^{2}\,f\!\left(r\right)}
\end{dmath*}
\begin{dmath*}[breakdepth={5}]
\left[P^{abcd}\right]^{1212}=\frac{f\!\left(r\right)}{2\,{r}^{2}}
\end{dmath*}
\bigskip\noindent\textbf{Política de lectura.} Este PDF/LaTeX es una vista. El archivo \texttt{results.json} conserva la representación intermedia canónica y la trazabilidad completa.
\end{document}
